\begin{center}
	\Large
	\textbf{KATA PENGANTAR}
\end{center}

\addcontentsline{toc}{chapter}{KATA PENGANTAR}

\vspace{2ex}

% Ubah paragraf-paragraf berikut dengan isi dari kata pengantar

Puji dan syukur ke hadirat Tuhan Yang Maha Esa atas segala rahmat dan karunia-Nya sehingga penulis dapat menyelesaikan penelitian yang berjudul ``Sistem Otonom Drone Holybro S500 untuk Pencarian Manusia Berbasis Kamera Monokular Tetap Menggunakan Raspberry Pi 5.''

Penelitian ini disusun dalam rangka memenuhi Tugas Akhir sebagai salah satu syarat memperoleh gelar Sarjana pada Departemen Teknik Komputer, Institut Teknologi Sepuluh Nopember. Oleh karena itu, penulis mengucapkan terima kasih kepada:

\begin{enumerate}[nolistsep]
	\item Bapak Dr. Arief Kurniawan, S.T., M.T., selaku Kepala Departemen Teknik Komputer, Fakultas Teknologi Elektro dan Informatika Cerdas, Institut Teknologi Sepuluh Nopember.
	\item Bapak \advisor, selaku Dosen Pembimbing I dan Bapak \coadvisor, selaku Dosen Pembimbing II yang telah memberikan bimbingan, arahan, serta dukungan kepada penulis selama proses penyusunan tugas akhir ini.
	\item Seluruh dosen Departemen Teknik Komputer yang telah memberikan ilmu, pengalaman, dan bimbingan kepada penulis selama menempuh pendidikan.
	\item Kedua orang tua dan seluruh keluarga yang senantiasa memberikan doa, kasih sayang, serta dukungan moral maupun material selama penulis menempuh pendidikan.
	\item Teman-teman Laboratorium Multimedia and Internet of Things (MIOT) yang telah memberikan bantuan, masukan, dan kerja sama selama proses pengerjaan tugas akhir ini.
	\item Teman-teman Departemen Teknik Komputer yang telah memberikan semangat, motivasi, dan kebersamaan selama masa perkuliahan.
	\item Seseorang yang senantiasa memberikan dukungan, semangat, motivasi, serta menjadi tempat berbagi selama proses penyusunan tugas akhir ini.
\end{enumerate}

Akhir kata, penulis berharap penelitian ini dapat memberikan manfaat bagi pengembangan ilmu pengetahuan maupun bagi pihak-pihak yang membutuhkan. Penulis menyadari bahwa tugas akhir ini masih memiliki berbagai kekurangan. Oleh karena itu, penulis mengharapkan kritik dan saran yang membangun demi penyempurnaan karya ini di masa mendatang.

\begin{flushright}
	\begin{tabular}[b]{c}
		\place{}, \MONTH{} \the\year{} \\
		\\
		\\
		\\
		\\
		\name{}
	\end{tabular}
\end{flushright}
